\hypertarget{c11-move-semantics}{%
\chapter{C++11 MOVE SEMANTICS}\label{c11-move-semantics}}

Move semantics is arguably the most significant performance feature in
C++11. It allows resources to be ``transferred'' (moved) from temporary
objects rather than copied.

\begin{center}\rule{0.5\linewidth}{0.5pt}\end{center}

\hypertarget{lvalues-and-rvalues}{%
\section{1. LVALUES AND RVALUES}\label{lvalues-and-rvalues}}

\hypertarget{lvalues}{%
\subsection{1.1 Lvalues}\label{lvalues}}

An \textbf{lvalue} (locator value) represents an object that occupies an
identifiable location in memory (has an address). - Example:
\passthrough{\lstinline!int x = 5;!} (\passthrough{\lstinline!x!} is
lvalue). - You can take its address: \passthrough{\lstinline!\&x!}.

\hypertarget{rvalues}{%
\subsection{1.2 Rvalues}\label{rvalues}}

An \textbf{rvalue} is everything else: temporary values, literals, or
results of expressions. - Example: \passthrough{\lstinline!5!},
\passthrough{\lstinline!x + 2!},
\passthrough{\lstinline!funcReturningVal()!}. - You cannot take its
address.

\begin{center}\rule{0.5\linewidth}{0.5pt}\end{center}

\hypertarget{rvalue-references}{%
\section{2. RVALUE REFERENCES}\label{rvalue-references}}

C++11 introduced the rvalue reference: \passthrough{\lstinline!T\&\&!}.
It binds \emph{only} to rvalues.

\begin{lstlisting}[language={C++}]
int x = 10;
int& lref = x;      // Lvalue ref binds to lvalue
// int&& rref = x;  // Error: cannot bind rvalue ref to lvalue

int&& rref2 = 20;   // OK: 20 is rvalue
\end{lstlisting}

\begin{center}\rule{0.5\linewidth}{0.5pt}\end{center}

\hypertarget{move-constructor-assignment}{%
\section{3. MOVE CONSTRUCTOR \&
ASSIGNMENT}\label{move-constructor-assignment}}

This allows a class to steal resources from a temporary object instead
of making a deep copy.

\hypertarget{deep-copy-the-old-way}{%
\subsection{3.1 Deep Copy (The Old Way)}\label{deep-copy-the-old-way}}

\begin{lstlisting}[language={C++}]
class Vector {
    int* data;
    size_t size;
public:
    // Copy Constructor
    Vector(const Vector& other) : size(other.size) {
        data = new int[size];
        std::copy(other.data, other.data + size, data);
    }
};
\end{lstlisting}

\hypertarget{move-constructor-the-c11-way}{%
\subsection{3.2 Move Constructor (The C++11
Way)}\label{move-constructor-the-c11-way}}

\begin{lstlisting}[language={C++}]
    // Move Constructor
    Vector(Vector&& other) noexcept : data(other.data), size(other.size) {
        // Steal the pointer
        other.data = nullptr; // Null out source
        other.size = 0;
    }
\end{lstlisting}

If \passthrough{\lstinline!other!} is a temporary, the compiler selects
the Move Constructor. This is O(1) instead of O(N).

\begin{center}\rule{0.5\linewidth}{0.5pt}\end{center}

\hypertarget{stdmove}{%
\section{4. STD::MOVE}\label{stdmove}}

\passthrough{\lstinline!std::move(x)!} does exactly one thing: it casts
\passthrough{\lstinline!x!} to an rvalue reference
(\passthrough{\lstinline!T\&\&!}). It essentially says, ``I am done with
this object, you can steal from it.''

\begin{lstlisting}[language={C++}]
Vector v1(100);
Vector v2 = std::move(v1); // Calls Move Constructor
// v1 is now empty (if implemented correctly)
\end{lstlisting}

\begin{center}\rule{0.5\linewidth}{0.5pt}\end{center}

\hypertarget{perfect-forwarding}{%
\section{5. PERFECT FORWARDING}\label{perfect-forwarding}}

Used in templates to preserve the value category (lvalue vs rvalue) of
arguments.

\hypertarget{universal-references-forwarding-references}{%
\subsection{5.1 Universal References (Forwarding
References)}\label{universal-references-forwarding-references}}

If \passthrough{\lstinline!T!} is a template parameter,
\passthrough{\lstinline!T\&\&!} is a \textbf{universal reference}, not
just an rvalue reference. It can bind to anything.

\hypertarget{stdforward}{%
\subsection{5.2 std::forward}\label{stdforward}}

\begin{lstlisting}[language={C++}]
template<typename T>
void wrapper(T&& arg) {
    func(std::forward<T>(arg));
}
\end{lstlisting}

\begin{itemize}
\tightlist
\item
  If \passthrough{\lstinline!wrapper!} is called with lvalue,
  \passthrough{\lstinline!arg!} is lvalue,
  \passthrough{\lstinline!forward!} keeps it lvalue.
\item
  If \passthrough{\lstinline!wrapper!} is called with rvalue,
  \passthrough{\lstinline!arg!} is lvalue (as a named variable), but
  \passthrough{\lstinline!forward!} casts it back to rvalue.
\end{itemize}

This enables \passthrough{\lstinline!emplace\_back!} to work
efficiently.
